% =====================================================================
%  Thème 2 — CORRIGÉ — Niveau FACILE
% =====================================================================
\documentclass[11pt,a4paper]{article}
\usepackage[utf8]{inputenc}
\usepackage[T1]{fontenc}
\usepackage{lmodern}
\usepackage[french]{babel}
\usepackage{amsmath,amssymb,mathtools}
\usepackage{icomma}
\usepackage{xcolor}
\usepackage[most]{tcolorbox}
\usepackage{enumitem}
\usepackage{array}
\usepackage{booktabs}
\usepackage{fancyhdr}
\usepackage[hidelinks]{hyperref}
\usepackage[a4paper,margin=2.2cm,headheight=26pt,headsep=10pt]{geometry}

\definecolor{primary}{RGB}{223,87,99}
\definecolor{primarydark}{RGB}{150,42,55}
\definecolor{excol}{RGB}{0,135,90}

\tcbset{boxbase/.style={enhanced,breakable,boxrule=0.9pt,arc=2.5pt,
   left=3mm,right=3mm,top=2.3mm,bottom=2.3mm,
   fonttitle=\bfseries,coltitle=white,toptitle=0.6mm,bottomtitle=0.6mm}}
\newtcolorbox[auto counter]{correction}[1][]{boxbase,colback=excol!5,colframe=excol,
  title={\textsc{Corrigé}~\thetcbcounter\ifx\relax#1\relax\relax\else\textmd{ --- \itshape #1}\fi}}

\pagestyle{fancy}
\fancyhf{}
\fancyhead[L]{\small\textcolor{primarydark}{\textbf{Fiche 6} — Les limites}}
\fancyhead[R]{\small\textcolor{primarydark}{\textsc{Thème 2 — Corrigé Facile}}}
\fancyfoot[C]{\small\thepage}
\renewcommand{\headrulewidth}{0.8pt}
\renewcommand{\headrule}{\hbox to\headwidth{\color{primary}\leaders\hrule height \headrulewidth\hfill}}

\begin{document}

\begin{center}
{\Large\bfseries\color{primarydark}Thème 2 — Règle $\dfrac{a}{0^{\pm}}$ \& asymptote verticale}\\[2pt]
{\large\color{excol}Corrigé — Niveau Facile}
\end{center}
\vspace{6pt}

\begin{correction}[appliquer la règle des signes]
Règle des signes ($0^{+}$ compte $+$, $0^{-}$ compte $-$) :
\begin{enumerate}[label=\alph*),leftmargin=2em,itemsep=3pt]
  \item $\dfrac{5}{0^{+}}=+\infty$ \quad b) $\dfrac{-3}{0^{+}}=-\infty$ \quad
        c) $\dfrac{2}{0^{-}}=-\infty$ \quad d) $\dfrac{-7}{0^{-}}=+\infty$ (moins par moins $=$ plus).
\end{enumerate}
\end{correction}

\begin{correction}[signe du dénominateur]
\begin{enumerate}[label=\alph*),leftmargin=2em,itemsep=3pt]
  \item $x\to 4^{+}$ : $x>4$ donc $x-4>0$, c.-à-d.\ $x-4\to 0^{+}$.
  \item $x\to 4^{-}$ : $x<4$ donc $x-4<0$, c.-à-d.\ $x-4\to 0^{-}$.
  \item $x\to 2^{+}$ : $x>2$ donc $2-x<0$, c.-à-d.\ $2-x\to 0^{-}$.
  \item $x\to 1^{-}$ : $(x-1)^2>0$ (un carré) donc $(x-1)^2\to 0^{+}$ (des deux côtés d'ailleurs).
\end{enumerate}
\end{correction}

\begin{correction}[limites latérales élémentaires]
\begin{enumerate}[label=\alph*),leftmargin=2em,itemsep=3pt]
  \item $\displaystyle\lim_{x\to 0^{+}}\frac{3}{x}=\frac{3}{0^{+}}=+\infty$ ; \quad
        $\displaystyle\lim_{x\to 0^{-}}\frac{3}{x}=\frac{3}{0^{-}}=-\infty$.
  \item $x\to 4^{+}$ : $x-4\to 0^{+}$, donc $\displaystyle\lim_{x\to 4^{+}}\frac{1}{x-4}=\frac{1}{0^{+}}=+\infty$ ;\\
        $x\to 4^{-}$ : $x-4\to 0^{-}$, donc $\displaystyle\lim_{x\to 4^{-}}\frac{1}{x-4}=\frac{1}{0^{-}}=-\infty$.
\end{enumerate}
\end{correction}

\begin{correction}[une fonction rationnelle simple]
$f(x)=\dfrac{7}{x+3}$.
\begin{enumerate}[label=\alph*),leftmargin=2em,itemsep=3pt]
  \item Le numérateur est la constante $7$ : il ne tend pas vers $0$, il est $>0$.
  \item $x\to -3^{+}$ : $x+3\to 0^{+}$, donc $\displaystyle\lim_{x\to -3^{+}} f(x)=\frac{7}{0^{+}}=+\infty$ ;\\
        $x\to -3^{-}$ : $x+3\to 0^{-}$, donc $\displaystyle\lim_{x\to -3^{-}} f(x)=\frac{7}{0^{-}}=-\infty$.
\end{enumerate}
\end{correction}

\begin{correction}[repérer une asymptote verticale]
On cherche le zéro du dénominateur (numérateur non nul en ce point) :
\begin{enumerate}[label=\alph*),leftmargin=2em,itemsep=3pt]
  \item $x-5=0\Rightarrow x_0=5$ (asymptote verticale $x=5$).
  \item $2x+6=0\Rightarrow x_0=-3$ (asymptote verticale $x=-3$).
  \item $x^2=0\Rightarrow x_0=0$ (asymptote verticale $x=0$ ; numérateur $5\neq 0$).
\end{enumerate}
\end{correction}

\end{document}
